\section{A Loader Behind the Reference Resolver}
\label{sec:ch10-the-loader-behind-the-resolver}

\index{DataLoader!behind a reference resolver|(}%
The fix is chapter~\ref{ch:data-without-n-plus-1}'s, in a different method.
Here is \code{SessionDataLoader.cs} in full, new in the Sessions service:

\begin{minted}{csharp}
using Microsoft.EntityFrameworkCore;

namespace Sessions;

public static class SessionDataLoader
{
    /// <summary>
    /// Every session one batch of representations asked for, keyed by
    /// identifier. The router sends this service a list rather than
    /// one key at a time, so the loader behind the reference resolver
    /// is what decides whether that list costs one statement or one
    /// statement each.
    /// </summary>
    // No projection here, and that is a decision rather than an
    // oversight. A QueryContext<Session> can be injected and does
    // narrow the SELECT, but the selector it builds fills every
    // property nobody asked for with a default, and Id is a property
    // nobody asks for. This dictionary is keyed on Id. Chapter 10
    // measures both halves and takes the columns rather than the risk.
    [DataLoader]
    public static async Task<Dictionary<int, Session>>
        GetSessionByIdAsync(
            IReadOnlyList<int> ids,
            ConferenceContext db,
            CancellationToken ct)
        => await db.Sessions
            .Where(s => ids.Contains(s.Id))
            .ToDictionaryAsync(s => s.Id, ct);
}
\end{minted}

\index{DataLoader!the same file, in a third service}%
That is \code{SpeakerDataLoader.cs} from
chapter~\ref{ch:data-without-n-plus-1} with a different table and a different
context. All three services in the graph now carry one, which is what
chapter~\ref{ch:second-third-subgraph}'s rule about sharing no code between
them costs and buys. The comment in the middle is
section~\ref{sec:ch10-what-the-loader-does-not-fix}'s subject and can wait.

\index{reference resolver!takes the loader, not the context}%
\code{SessionType.cs} changes by one parameter, one line of body and the doc
comment, which now says what the loader is for instead of naming this chapter:

\begin{minted}[highlightlines={2-8,14,29,30}]{csharp}
    /// <summary>
    /// Answers one representation from the entity route. Two things
    /// differ from the node resolver above. The key arrives as the
    /// undecoded node id string, so this method decodes it. And the
    /// router hands this route a whole list rather than one key at a
    /// time, so the lookup goes through a DataLoader where the one
    /// above projects instead; SessionDataLoader says why it cannot
    /// do both.
    /// </summary>
    [GraphQLIgnore]
    public static async Task<Session?> ResolveSessionReferenceAsync(
        string id,
        INodeIdSerializer serializer,
        ISessionByIdDataLoader sessionById,
        CancellationToken ct)
    {
        var nodeId = serializer.Parse(id, typeof(int));

        // Nothing has checked what is inside this string, and the
        // router is not the only caller the entity route has. A
        // Speaker's node id decodes to a perfectly good integer;
        // without this line it would answer with the session that
        // happens to carry it.
        if (nodeId.TypeName != nameof(Session))
        {
            return null;
        }

        return await sessionById.LoadAsync(
            (int)nodeId.InternalId!, ct);
    }
\end{minted}

\index{DbContext@\code{DbContext}!the resolver stops knowing about one}%
The \code{ConferenceContext} parameter is gone and the method no longer knows
there is a database, which is the same thing that happened to
\code{SessionType.GetSpeakerAsync} in
chapter~\ref{ch:data-without-n-plus-1} and is worth as much here. What the
method does now is say which session it wants and check that it was asked a
question about a session.

\index{statement count!three representations, one statement}%
Send the three representations again, on the Sessions service's own port so
that no router is involved:

\begin{minted}{text}
-- statement 1
SELECT "s"."Id", "s"."Abstract", "s"."DurationMinutes",
       "s"."SpeakerId", "s"."StartsAt", "s"."Title"
FROM "Sessions" AS "s"
WHERE "s"."Id" IN (@ids1, @ids2, @ids3)
\end{minted}

\index{statement count!the whole graph is back to two}%
One. And through the router, the page that opened this chapter is back to the
two statements chapter~\ref{ch:second-third-subgraph} measured, one in each
service, with the Speakers half being the loader that was there all along.

\subsection{Four Things That Have to Survive}

\index{entities@\code{\_entities}!batching must not disturb the contract}%
A fix that quietly broke the entity route's contract would be worse than the
N+1, and three of the four things below are that contract. A dictionary is a
very different object from a row, and what puts each row back where its
representation was is somebody else's code, so none of the four is safe to
assume.

\index{entities@\code{\_entities}!order survives batching}%
Order is the first requirement of the specification and the one a batch is most
likely to lose, because a dictionary has no order at all. Ask for three
sessions in an order that is not key order:

\begin{minted}[fontsize=\footnotesize,breakanywhere]{json}
{"r":[{"__typename":"Session","id":"U2Vzc2lvbjoz"},
      {"__typename":"Session","id":"U2Vzc2lvbjox"},
      {"__typename":"Session","id":"U2Vzc2lvbjoy"}]}
\end{minted}

\begin{minted}[fontsize=\footnotesize,breakanywhere]{json}
{"data":{"_entities":[
  {"__typename":"Session","title":"Paging Without Offsets"},
  {"__typename":"Session","title":"Schemas That Outlive Their Authors"},
  {"__typename":"Session","title":"Reading a Query Plan"}]}}
\end{minted}

The order asked for, not the order the database found them in, and still one
statement. The engine keeps each representation's position and fills it from the
dictionary the loader returned. Chapter~\ref{ch:representation-order} is the
chapter about what a response looks like when that does not hold, and it is
worth knowing before you read it that a DataLoader is not how it breaks.

\index{entities@\code{\_entities}!a missing row is still a null}%
A key that matches no row is the second. Chapter~\ref{ch:first-subgraph}
established that it is a null in the list and not an error, which is what the
specification permits. Put one between two that do match, by sending session
one, session ninety-nine and session two:

\begin{minted}[fontsize=\footnotesize,breakanywhere]{json}
{"data":{"_entities":[
  {"__typename":"Session","title":"Schemas That Outlive Their Authors"},
  null,
  {"__typename":"Session","title":"Reading a Query Plan"}]}}
\end{minted}

Still a null, still no error, still one statement.
Chapter~\ref{ch:data-without-n-plus-1} noted that the dictionary a loader
returns does not have to contain every key it was given and that a key left out
resolves to null, and this is where that default earns its keep: the key the
loader leaves out is exactly the entry the specification wants nulled.

\index{DataLoader!de-duplicates for every caller, not just the router}%
De-duplication is the third, and it now reaches callers that are not the
router. Section~\ref{sec:ch10-three-representations-not-four} sent the same
representation three times and got three statements. Now it is one, carrying
one key, and three identical answers come back. The router had been
de-duplicating on its own behalf; the loader does it for anybody.

\index{reference resolver!the guard runs before the batch}%
The guard is the fourth. Chapter~\ref{ch:first-subgraph} put a type-name check
in front of the key, and showed what the entity route answers without one.
Batching could plausibly have moved it: if keys were collected first and
checked later, a rejected one would still cost a lookup. It does not. Send a
\code{Speaker} node id under a \code{Session} type name in the middle of two
good ones and the statement reads:

\begin{minted}{text}
WHERE "s"."Id" IN (@ids1, @ids2)
\end{minted}

Two keys for three representations, and a null in the middle position. The
resolver runs per representation, decides per representation, and only the ones
that get as far as \code{LoadAsync} are in the batch. The guard still costs
nothing.

\subsection{The Third Shape, Which Needs None of This}

\index{Ratings service!a reference resolver with no table behind it}%
Before you go and put a loader behind every reference resolver you own, the
Ratings service is the counter-example and it is in this graph.
Chapter~\ref{ch:second-third-subgraph} printed that service's reference
resolver in full, and the shape of the method is the whole argument: it is not
async, it takes no \code{DbContext} and no loader, and it issues no statement.

\index{Ratings service!answers out of the representation}%
It builds a \code{Session}
out of the representation it was handed, which is why
chapter~\ref{ch:second-third-subgraph} measured \code{feedbackUrl} at zero
statements in that service. A reference resolver that touches no table has no
N+1 to have, and a DataLoader in front of it would be a cache with nothing
behind it.

\index{DataLoader!the request cache, not the batch}%
The Ratings service does use a loader, and for the other of the two things a
loader does. \code{averageScore} and \code{ratingCount} are ordinary field
resolvers on the entity, and both call \code{LoadAsync} with the same key. Four
sessions, two fields, eight loads, one statement. The batch explains the first
factor and the request-scoped cache explains the second, and the two have been
separate since the original library, whose README says both: it
\enquote{will coalesce all individual loads which occur within a single frame of
execution} and, of the second, that
\enquote{subsequent calls to .load() with the same key will result in that key
not appearing in the keys provided to your batch
function}~\autocite{dataloaderjs}. Keep them apart in your head, because only
the first is what this chapter is about.

\index{sources!who publishes this advice, and who does not}%
So the rule I would give you is narrower than the one Apollo gives. Apollo's
documentation says a reference resolver
\enquote{might be called many times to resolve a single query} and that it is
\enquote{crucial that reference resolvers account for \enquote{N+1} issues
(typically via data loaders)}, and then generalizes to using one
\enquote{in every resolver}~\autocite{apolloentities, apollonplusone}. The
narrow half is right and I have just measured it. The general half I would not
write: put a loader where a list of keys meets a store, and leave it out where
there is no store.
\index{DataLoader!behind a reference resolver|)}%
